% Resumo em língua vernácula
\setlength{\absparsep}{18pt}
\begin{resumo}
As vigas mistas de aço e concreto biapoiadas, associadas a lajes com fôrma de aço incorporada (\emph{steel deck}) ou a lajes maciças e executadas sem escoramento, são a solução predominante para pisos de edifícios comerciais no Brasil. O seu dimensionamento é multifásico: na concretagem (Fase 1) o perfil de aço resiste isoladamente ao peso próprio, ao concreto fresco e à sobrecarga de construção; após a cura (Fase 2) a seção trabalha como mista, com a transferência do fluxo de cisalhamento longitudinal feita por conectores. A verificação envolve estados-limites últimos (flexão, força cortante e grau de interação da conexão) e de serviço (deslocamentos com fluência e retração e, a partir da edição de 2024 da ABNT NBR 8800, vibração do piso pelo caminhar de pessoas, no Anexo I). Este trabalho desenvolve uma ferramenta computacional de código aberto, escrita em C++20 com interface gráfica em Dear ImGui e Raylib, que automatiza esse encadeamento para as edições de 2008 e de 2024 da norma, mantendo as duas implementações separadas para permitir a comparação direta. O trabalho apresenta a fundamentação do Anexo L, a formalização dos algoritmos e a estratégia de verificação do software por testes automatizados. Para validação, adotou-se o exemplo de um piso de escritórios com viga W460\ensuremath{\times}52,0, \emph{steel deck} MF-75 e vão de 9,0 m, resolvido em uma planilha de referência. O software reproduziu as 25 grandezas comparadas (esforços, resistências, grau de interação, inércias efetivas, flechas e parâmetros de vibração) com diferença relativa máxima de $2{,}4 \times 10^{-15}$; o procedimento de vibração foi conferido também contra um segundo exemplo independente. Os estudos paramétricos mostraram que, no caso analisado, a edição de 2024 reduz o grau de interação mínimo de 0,519 para 0,280, mas, em lajes maciças, exige cerca de 30\% mais conectores para a interação completa, em razão do novo coeficiente $R_p = 0{,}75$; mostraram também que o critério indireto de vibração da edição de 2008 reprova a viga do estudo de caso, enquanto a avaliação pelo Anexo I a aprova pelo critério da aceleração de pico. A arquitetura modular adotada serve de base para a extensão futura a vigas mistas protendidas.

\textbf{Palavras-chave:} Vigas mistas aço-concreto; ABNT NBR 8800:2024; Ferramenta computacional; Construção não escorada; \emph{Steel deck}; Grau de interação; Vibração de pisos.
\end{resumo}
